\chapter{Elliptic Operators on Compact Manifolds}\label{elliptic-operators}
\chapterstatus{complete}

\section{Introduction}\label{elliptic-operators:section-introduction}

This chapter contains the analysis behind the Hodge theorem. An elliptic differential operator on a compact manifold is Fredholm
between Sobolev spaces (Theorem~\ref{elliptic-operators:theorem-fredholm}), its kernel consists of smooth sections, and a
self-adjoint elliptic operator has a discrete spectrum with smooth eigensections
(Theorem~\ref{elliptic-operators:theorem-spectral}). Applied to the Hodge Laplacian these statements give the decomposition of the
space of forms and the existence of harmonic representatives (Chapter~\ref{hodge-theorem}). References:
\cite[Chapter 6]{warner-foundations}, \cite[Chapter 4]{wells-differential-analysis}, \cite{hormander-pdo1}.

\noindent\textbf{Prerequisites.} Chapter~\ref{riemannian-geometry} (Riemannian Geometry) and Chapter~\ref{distributions}
(Distributions and Sobolev Spaces).

Throughout, $M$ is a compact smooth manifold, and $E$, $F$ are complex vector bundles on $M$ with Hermitian metrics; $M$ carries a
fixed Riemannian metric and, when needed, an orientation, so that sections can be integrated.

\section{Differential operators and symbols}\label{elliptic-operators:section-symbols}

\begin{definition}\label{elliptic-operators:definition-differential-operator}
A linear map $P \colon \Gamma(M, E) \to \Gamma(M, F)$ is a \emph{differential operator of order at most $m$} if in local coordinates and local frames it is given by
$Pu = \sum_{|\alpha| \leq m}A_\alpha(x)\,\partial^\alpha u$ with smooth matrix-valued coefficients $A_\alpha$. Its \emph{principal symbol} is
\[
\sigma_m(P)(x, \xi) = \sum_{|\alpha| = m}A_\alpha(x)\,(i\xi)^\alpha \in \Hom(E_x, F_x), \qquad \xi \in T^*_xM,
\]
and $P$ is \emph{elliptic} if $\sigma_m(P)(x, \xi)$ is invertible for all $x$ and all $\xi \neq 0$.
\end{definition}

\begin{proposition}\label{elliptic-operators:proposition-symbol}
The principal symbol is well defined: it does not depend on the coordinates and frames, and it is characterised by
\[
\sigma_m(P)(x, d_xf)\,u(x) = \frac{i^m}{m!}\,P\big((f - f(x))^mu\big)(x)
\]
for $f \in C^\infty(M)$ and $u \in \Gamma(M, E)$. Symbols are multiplicative, $\sigma_{m+m'}(P'P) = \sigma_{m'}(P')\sigma_m(P)$, so a composite of elliptic operators is elliptic, and the symbol of the formal
adjoint is the adjoint of the symbol.
\end{proposition}

\begin{proof}
Expanding $P((f - f(x))^mu)$ by the Leibniz rule, all terms in which fewer than $m$ derivatives fall on $(f - f(x))^m$ vanish at $x$, and the remaining term is
$\sum_{|\alpha| = m}A_\alpha(x)\,\partial^\alpha\big((f - f(x))^m\big)(x)u(x) = m!\sum_{|\alpha| = m}A_\alpha(x)(d_xf)^\alpha u(x)$, which gives the formula. The right-hand side is defined without reference to
coordinates, so the symbol is invariantly defined, and multiplicativity follows from the same formula applied twice.
\end{proof}

\begin{example}\label{elliptic-operators:example-elliptic}
\begin{enumerate}
\item The Hodge Laplacian $\Delta = d\delta + \delta d$ on $\Omega^k(M)$ has principal symbol $|\xi|^2\,\id$ (compute in normal coordinates, where the first-order terms of the metric vanish); it is elliptic.
\item The operator $d + \delta \colon \Omega^{\mathrm{even}} \to \Omega^{\mathrm{odd}}$ has symbol $i(\xi \wedge - \iota_\xi)$, whose square is $-|\xi|^2$, so it is invertible for $\xi \neq 0$ and $d + \delta$ is
elliptic. Its index is the Euler characteristic (Chapter~\ref{hodge-theorem}).
\item On a complex manifold, $\bar\partial + \bar\partial^*$ is elliptic (Chapter~\ref{dolbeault}); the Cauchy--Riemann operator $\bar\partial$ itself is elliptic in one variable.
\item Not every operator is elliptic: the wave operator $\partial_t^2 - \Delta$ and the heat operator $\partial_t - \Delta$ are not.
\end{enumerate}
\end{example}

\begin{proposition}[Fundamental solution of the Laplacian]\label{elliptic-operators:proposition-fundamental-solution}
On $\RR^3$, the locally integrable function $E(x) = -\frac{1}{4\pi|x|}$ satisfies $\sum_i\partial_i^2E = \delta_0$ in $\mathcal{D}'(\RR^3)$. Consequently, for $f \in \mathcal{D}(\RR^3)$ the convolution $u = E * f$ solves $\sum_i\partial_i^2u = f$.
\end{proposition}

\begin{proof}
$E$ is locally integrable, as $\int_{|x| \leq 1}|x|^{-1}\,dx = 4\pi\int_0^1r\,dr$ in spherical coordinates (Theorem~\ref{measure-theory:theorem-change-of-variables}). Away from $0$, a direct computation gives $\Delta E = 0$, where $\Delta = \sum_i\partial_i^2$ here. Let
$\varphi \in \mathcal{D}(\RR^3)$ and $A_\varepsilon = \{\varepsilon \leq |x| \leq R\}$ with $R$ large. By the divergence theorem (Example~\ref{differential-forms:example-classical}) applied to $E\nabla\varphi - \varphi\nabla E$ on $A_\varepsilon$ (Green's
identity), and since $\varphi$ vanishes near $|x| = R$,
\[
\int_{A_\varepsilon}E\,\Delta\varphi = \int_{A_\varepsilon}(E\,\Delta\varphi - \varphi\,\Delta E) = -\int_{|x| = \varepsilon}\big(E\,\partial_r\varphi - \varphi\,\partial_rE\big)\,dA,
\]
where $\partial_r$ is the radial derivative, the sign coming from the inner boundary sphere having outward normal $-x/|x|$. On $|x| = \varepsilon$, $E = -1/(4\pi\varepsilon)$ and $\partial_rE = 1/(4\pi\varepsilon^2)$; the area is $4\pi\varepsilon^2$. So the first term is
at most $\varepsilon\sup|\nabla\varphi| \to 0$, and the second is $\frac{1}{4\pi\varepsilon^2}\int_{|x| = \varepsilon}\varphi\,dA \to \varphi(0)$ by continuity. Letting $\varepsilon \to 0$ (dominated convergence on the left),
$(\Delta E)(\varphi) = \int E\,\Delta\varphi = \varphi(0)$. For the last statement, $\Delta(E * f) = E * \Delta f$ by differentiation under the integral sign, and $(E * \Delta f)(x) = \int E(y)\,\Delta_y f(x - y)\,dy = f(x)$ by the first part applied
to $\varphi(y) = f(x - y)$.
\end{proof}

\section{Sobolev spaces of sections}\label{elliptic-operators:section-sobolev-sections}

\begin{definition}\label{elliptic-operators:definition-sobolev-sections}
Fix a finite atlas of $M$ with trivialisations of $E$ and a subordinate partition of unity $(\psi_i)$. For $s \in \RR$ and $u \in \Gamma(M, E)$ set
$\|u\|_s^2 = \sum_i\|\psi_iu\|_{H^s(\RR^n)}^2$, using the coordinates and frames to view $\psi_iu$ as a compactly supported $\CC^r$-valued function. The completion is the Sobolev space $H^s(E)$ of sections.
\end{definition}

\begin{proposition}\label{elliptic-operators:proposition-sobolev-sections}
The space $H^s(E)$ and the equivalence class of its norm do not depend on the choices. Moreover $H^0(E) = L^2(E)$ with the inner product $\int_M\langle u, v\rangle\,\mathrm{vol}$; the inclusions
$H^s(E) \to H^t(E)$ for $s > t$ are compact; $\bigcap_sH^s(E) = \Gamma(M, E)$; and a differential operator of order $m$ extends to a bounded map $H^{s+m}(E) \to H^s(F)$ for every $s$.
\end{proposition}

\begin{proof}
Changing charts and frames changes $\psi_iu$ by composition with a diffeomorphism and multiplication by a smooth matrix, which are bounded operations on $H^s$ of compactly supported functions
\cite[Chapter 6]{warner-foundations}; hence the norms are equivalent. For $s = 0$ the statement is Plancherel's theorem plus the fact that the two ways of integrating differ by a bounded density. Compactness
follows from the Rellich theorem (Theorem~\ref{distributions:theorem-rellich}, in the form of Counterexample~\ref{distributions:counterexample-rellich-rn} for compactly supported distributions) applied to
each $\psi_iu$. The identification of $\bigcap_sH^s$ with smooth sections is the Sobolev embedding theorem (Theorem~\ref{distributions:theorem-sobolev-embedding}), and boundedness of differential operators
follows because differentiation lowers $s$ by the order and multiplication by smooth functions is bounded.
\end{proof}

\section{Elliptic estimates and regularity}\label{elliptic-operators:section-estimates}

\begin{theorem}\label{elliptic-operators:theorem-elliptic-estimate}
Let $P$ be an elliptic differential operator of order $m$ from $E$ to $F$ over the compact manifold $M$. For every $s$ there is $C_s$ with
\[
\|u\|_{s+m} \leq C_s\left(\|Pu\|_s + \|u\|_s\right) \qquad \text{for all } u \in H^{s+m}(E),
\]
and (\emph{elliptic regularity}) if $u \in H^t(E)$ for some $t$ and $Pu \in H^s(F)$, then $u \in H^{s+m}(E)$; in particular $Pu$ smooth implies $u$ smooth.
\end{theorem}

\begin{proof}[Sketch of proof]
Locally one freezes the coefficients: near a point $x_0$, $P$ differs from the constant-coefficient elliptic operator $P_0 = \sum_{|\alpha| = m}A_\alpha(x_0)\partial^\alpha$ by an operator whose top-order
coefficients are small on a small neighbourhood and whose remaining terms have order $< m$. The constant-coefficient estimate
(Theorem~\ref{distributions:theorem-elliptic-estimate}, transported to $\RR^n$ by a cut-off and to the torus) then gives
$\|u\|_{s+m} \leq C(\|Pu\|_s + \|u\|_{s+m-1})$ for $u$ supported in that neighbourhood, and the error term is absorbed using the interpolation inequality
$\|u\|_{s+m-1} \leq \varepsilon\|u\|_{s+m} + C_\varepsilon\|u\|_s$, which follows from
$\langle\xi\rangle^{s+m-1} \leq \varepsilon\langle\xi\rangle^{s+m} + C_\varepsilon\langle\xi\rangle^s$. Summing over a finite cover with a partition of unity gives the global estimate. Regularity is proved by the same
estimates applied to difference quotients, or by constructing a \emph{parametrix}: an operator $Q$ with $QP = I + S$ and $PQ = I + S'$ where $S$, $S'$ are smoothing operators; then $u = Q(Pu) - Su$ is as
regular as $Pu$ allows. The omitted verifications are the construction of the parametrix and the coordinate invariance of the estimates; see \cite[Chapter 6]{warner-foundations},
\cite[Chapter 4]{wells-differential-analysis} or \cite[Chapters 17--18]{hormander-pdo1}.
\end{proof}

\section{Fredholm theory and the spectral decomposition}\label{elliptic-operators:section-fredholm}

\begin{definition}\label{elliptic-operators:definition-formal-adjoint}
The \emph{formal adjoint} $P^*$ of a differential operator $P$ is the differential operator of the same order characterised by $\int_M\langle Pu, v\rangle = \int_M\langle u, P^*v\rangle$ for smooth sections;
it exists, is unique, and has symbol $\sigma_m(P^*)(x, \xi) = (-1)^m\sigma_m(P)(x, \xi)^*$, so $P^*$ is elliptic if $P$ is. The operator $P$ is \emph{formally self-adjoint} if $P = P^*$; the Hodge Laplacian
is (Proposition~\ref{riemannian-geometry:proposition-adjoint}).
\end{definition}

\begin{theorem}\label{elliptic-operators:theorem-fredholm}
Let $P$ be elliptic of order $m$ on a compact manifold. Then:
\begin{enumerate}
\item $\ker P = \{u \in \Gamma(M, E) : Pu = 0\}$ is finite-dimensional and consists of smooth sections, and the same holds for $\ker P^*$;
\item $P \colon H^{s+m}(E) \to H^s(F)$ is Fredholm for every $s$, with kernel and cokernel independent of $s$;
\item there is an orthogonal decomposition $L^2(F) = \operatorname{im}(P) \oplus \ker P^*$, and $Pu = f$ is solvable with $u$ smooth if and only if $f$ is smooth and orthogonal to $\ker P^*$.
\end{enumerate}
\end{theorem}

\begin{proof}
(1) Elements of the kernel are smooth by elliptic regularity. On $\ker P$ the elliptic estimate gives $\|u\|_{s+m} \leq C_s\|u\|_s$, so the $H^{s}$ and $H^{s+m}$ norms are equivalent there; since the inclusion
$H^{s+m} \to H^s$ is compact (Proposition~\ref{elliptic-operators:proposition-sobolev-sections}), the unit ball of $\ker P$ is compact, and $\ker P$ is finite-dimensional
(Proposition~\ref{functional-analysis:proposition-banach-examples}(3)).

(2) Closed range: let $K \subseteq H^{s+m}$ be the orthogonal complement of $\ker P$. We claim $\|u\|_{s+m} \leq C\|Pu\|_s$ on $K$. If not, there are $u_j \in K$ with $\|u_j\|_{s+m} = 1$ and $Pu_j \to 0$; by
compactness of the inclusion, after passing to a subsequence $u_j \to u$ in $H^s$, and the elliptic estimate gives
$\|u_j - u_l\|_{s+m} \leq C(\|P(u_j - u_l)\|_s + \|u_j - u_l\|_s) \to 0$, so $u_j \to u$ in $H^{s+m}$ with $\|u\|_{s+m} = 1$, $Pu = 0$ and $u \in K$, a contradiction. Hence $P$ has closed range, and with (1)
applied to $P$ and $P^*$ it is Fredholm; the kernel is independent of $s$ by regularity, and so is the cokernel by (3).

(3) The operator $P$ with domain $H^m(E)$ and its formal adjoint satisfy: $f \in L^2(F)$ is orthogonal to $\operatorname{im}(P)$ if and only if $\int\langle Pu, f\rangle = 0$ for all smooth $u$, that is,
$P^*f = 0$ in the sense of distributions, which by elliptic regularity means $f \in \ker P^*$ and $f$ smooth. Since the range is closed, $L^2 = \overline{\operatorname{im}P} \oplus \ker P^* = \operatorname{im}P \oplus \ker P^*$, where
$\operatorname{im}P$ is taken on $H^m$; and if $f$ is smooth and orthogonal to $\ker P^*$, a solution $u \in H^m$ exists and is smooth by regularity.
\end{proof}

\begin{theorem}\label{elliptic-operators:theorem-spectral}
Let $P$ be a formally self-adjoint elliptic operator of order $m > 0$ on a compact manifold, with $P \geq 0$ in the sense that $\int\langle Pu, u\rangle \geq 0$. Then $L^2(E)$ has a complete orthonormal
system of smooth eigensections of $P$, the eigenvalues $0 \leq \lambda_1 \leq \lambda_2 \leq \cdots$ are real, have finite multiplicity, and tend to $+\infty$. Every $u \in L^2(E)$ is the $L^2$-limit of its
eigenfunction expansion, which converges in $\Gamma(M, E)$ when $u$ is smooth.
\end{theorem}

\begin{proof}
By Theorem~\ref{elliptic-operators:theorem-fredholm} applied to $P + I$, which is elliptic and injective on smooth sections (as $\int\langle (P+I)u, u\rangle \geq \|u\|_{L^2}^2$), the map
$P + I \colon H^m(E) \to L^2(E)$ is bijective, with bounded inverse $G$ by the open mapping theorem (Theorem~\ref{functional-analysis:theorem-baire-consequences}). Composing with the compact inclusion
$H^m \to L^2$ shows that $G \colon L^2 \to L^2$ is compact, and it is self-adjoint and positive because $P + I$ is. By the spectral theorem for compact self-adjoint operators
(Theorem~\ref{functional-analysis:theorem-spectral-compact}) there is an orthonormal basis of $L^2(E)$ of eigenvectors of $G$, with eigenvalues $\mu_k > 0$ tending to $0$ (they are nonzero since $G$ is
injective). Then $Pu_k = (\mu_k^{-1} - 1)u_k$, so the $u_k$ are eigensections of $P$ with eigenvalues $\lambda_k = \mu_k^{-1} - 1 \to \infty$, smooth by elliptic regularity, and $\lambda_k \geq 0$ by positivity.
Finite multiplicity is finite-dimensionality of the eigenspaces, which are kernels of the elliptic operators $P - \lambda_k$. For smooth $u$, the coefficients $\langle u, u_k\rangle$ decay faster than any power
of $\lambda_k$, because $\lambda_k^N\langle u, u_k\rangle = \langle P^Nu, u_k\rangle$ is bounded; with the Sobolev embedding this gives convergence in every $H^s$, hence in $\Gamma(M, E)$.
\end{proof}

\begin{remark}\label{elliptic-operators:remark-index}
The index $\operatorname{ind}P = \dim\ker P - \dim\ker P^*$ depends only on the homotopy class of the principal symbol, and the Atiyah--Singer index theorem computes it as a characteristic number of the
symbol; for $d + \delta$ this gives the Gauss--Bonnet theorem, and for the Dolbeault operator the Riemann--Roch--Hirzebruch theorem (Chapter~\ref{kodaira}). See \cite{lawson-michelsohn}.
\end{remark}

\section{Exercises}\label{elliptic-operators:section-exercises}

\begin{exercise}\label{elliptic-operators:exercise-symbol-laplacian}
Compute the principal symbol of $d$, of $\delta$ and of $\Delta$ on forms, and verify that $\Delta$ is elliptic while $d$ is not (for $n \geq 2$).
\end{exercise}

\begin{solution}
$\sigma_1(d)(x, \xi)\alpha = i\xi \wedge \alpha$ and $\sigma_1(\delta)(x, \xi)\alpha = -i\iota_\xi\alpha$, so
$\sigma_2(\Delta)(x, \xi) = \xi \wedge \iota_\xi + \iota_\xi(\xi \wedge -) = |\xi|^2\id$ by the algebraic identity of
Proposition~\ref{multilinear-algebra:proposition-interior-antiderivation}. The symbol of $d$ is not injective on $\Lambda^k$ for $k \geq 1$, since $\xi \wedge \xi = 0$.
\end{solution}

\begin{exercise}\label{elliptic-operators:exercise-eigenvalues-torus}
Compute the spectrum of the Laplacian on functions on the flat torus $\RR^n/\ZZ^n$ and verify Theorem~\ref{elliptic-operators:theorem-spectral} directly.
\end{exercise}

\begin{solution}
The functions $e^{2\pi ik\cdot x}$, $k \in \ZZ^n$, are eigenfunctions with eigenvalues $4\pi^2|k|^2$, and they form an orthonormal basis of $L^2$
(Example~\ref{functional-analysis:example-fourier}); multiplicities are the numbers of lattice points on spheres, and the eigenvalues tend to infinity.
\end{solution}

\begin{exercise}\label{elliptic-operators:exercise-index-self-adjoint}
Show that a formally self-adjoint elliptic operator on a compact manifold has index zero, and that the index of an elliptic operator changes sign when $P$ is replaced by $P^*$.
\end{exercise}

\begin{exercise}\label{elliptic-operators:exercise-no-solutions}
Give an example of a compact manifold and a smooth function $f$ for which $\Delta u = f$ has no solution, and characterise the $f$ for which it does.
\end{exercise}

\begin{solution}
On a compact connected oriented $M$, $\ker\Delta$ on functions consists of the constants, so by Theorem~\ref{elliptic-operators:theorem-fredholm} the equation is solvable exactly when
$\int_Mf\,\mathrm{vol} = 0$; any $f$ with nonzero integral, for instance $f = 1$, gives no solution.
\end{solution}
